%% AMS-LaTeX Created with the Wolfram Language : www.wolfram.com

\documentclass[12pt]{article}

\usepackage[letterpaper,margin=1in]{geometry}
\usepackage{amsmath, amssymb, graphics, setspace}
\usepackage{graphicx} 
\graphicspath{ {./figs/} }
\usepackage{tcolorbox}
\usepackage{bm}
\usepackage[framed,numbered,autolinebreaks,useliterate]{mcode}
\newcommand{\mathsym}[1]{{}}
\newcommand{\unicode}[1]{{}}

\usepackage{hyperref}
\hypersetup{
    colorlinks=true,
    linkcolor=blue,
    filecolor=magenta,      
    urlcolor=cyan,
    pdftitle={Overleaf Example},
    pdfpagemode=FullScreen,
}

\def\mathbi#1{\textbf{\em #1}}

\begin{document}
\doublespacing	

\title{CE 401/5501 Homework 10\\
\vspace{2in}
\textbf{Name: \rule{2in}{0.15mm}}}
\author{}
\maketitle

\newcommand*{\I}{\imath} % imaginary unit i

\newpage
\noindent Please review lecture notes:
\begin{itemize}
    \item \href{https://umsystem.instructure.com/courses/280496/files/33698586?module_item_id=9919505}{Topic 8}: Generative Learning
    \item \href{https://umsystem.instructure.com/courses/280496/files/33699788?module_item_id=9919650}{Topic 10}: Generalization and model selection 
\end{itemize}

complete the following concept understanding questions:

\noindent \pmb{Q1} \\
Given two random variables $x$ and $c$ that $x$ is feature with $x\in \mathbb{R}$ and $c \in \{A1, A2, A3\}$ is a categorical variable. The following priors are known:
\begin{align*}
    P(A1) & = 0.1 \\
    P(A2) & = 0.3 \\
    P(A3) & = 0.6
\end{align*}

\begin{itemize}
    \item Given these priors, attempt to write out all the possible inference equations, i.e., the maximum a posteriori (MAP) expression. 

For example, we have the first one:
\[
P(A1 | x) \propto p(x|A1) \times 0.1 = 0.1 p(x|A1)
\]
Write out the other two.

\item Suppose we know the true distributions for the likelihood functions, which all follow an Exponential distribution 
\[
p(x|c) = \lambda e^{-\lambda x}
\]
with if $c = A1$, $\lambda = 0.5$; if $c = A2$, $\lambda = 1.0$; and if $c = A3$, $\lambda = 1.5$, plot all the resulting MAP expressions in one plot. 

\item Then suppose we have three measurements: $x = 0.5$, $x = 1$, and $x = 2$, can you decide the class labels for the three measurements?
\end{itemize}


\newpage
\noindent \pmb{Q2} \\
In this problem, we will do the practice example again as we did in class with some modifications (see Topic 8 notes).

Again we monitor the direct vertical loads from passing vehicles on a bridge. Based on engineering judgment, we assume that the loads follow different uniform distributions for two vehicle types:

\begin{itemize}
    \item \textbf{Car (class 0):} The load follows a uniform distribution 
    \[
    X \sim U(a,\, 2a+1).
    \]
    \item \textbf{Pickup Truck (class 1):} The load follows a uniform distribution 
    \[
    X \sim U(2a,\, 3a+2).
    \]
\end{itemize}

We observe the following data:
\[
\{(2, 0),\, (3, 0),\, (4, 0/1),\, (4, 1),\, (5, 1)\},
\]
where the first entry is the measured load and the second entry indicates the vehicle type (with 0 for a car and 1 for a pickup truck).

Yet, there is a `glitch' in that at the third data point, the observer (maybe through a camera) does not recognize whether the vehicle is a car or a truck. Given such uncertainty and since we have considered the case that it is a car ('0'), let's assume that it could be a pickup truck (`1').

\begin{itemize}
    \item Follow the class example process, do it again, and determine the resulting total likelihood and the estimate to the parameter $a$.
    \item With the estimate of $a$, calculate the total likelihood from the two cases, namely assuming it is a car and it is a pickup truck. By comparing the total likelihood values, which scenario is more reliable?
    \item In class, we consider, which is a strong assumption, equal priors, namely 
    \begin{align*}
    P(`car') & = P(0) = 0.5 \\
    P(`pickup') &= P(1) = 0.5
    \end{align*}

If we keep this strong assumption, does it change Baye's error rate?

What if we consider the reality (the evidence of five observations we have), considering the two scenarios among the five observations:
\begin{itemize}
\item if it is a car:
  \begin{align*}
    P(`car') & = P(0) = \frac{3}{5} \\
    P(`pickup') &= P(1) = \frac{2}{5}
    \end{align*}
\item if it is a pickup
  \begin{align*}
    P(`car') & = P(0) = \frac{2}{5} \\
    P(`pickup') &= P(1) = \frac{3}{5}
    \end{align*}
\end{itemize}
Does it modify the underlying Baye's error rate?
    
\end{itemize}

\newpage
\noindent \pmb{Q3} \\
Regarding the generalization performance of a classification model (Topic 10), denoted as 
\[
m(c|\bm{x}, \bm{w}, \alpha)
\]
which outputs the class label given an input $\bm{x}$, where $\bm{w}$ is the model weight vector and $\alpha$ is a hyperparameter (e.g., the $C$ parameter in an SVM or a threshold in an MLP (at the last layer to threshold the sigmoid output, in general, it is 0.5 but it could be selected differently). 

Select from the following statements which one is very likely not correct.

\begin{enumerate}
    \item A model that underfits the training data typically exhibits high bias and low variance, leading to poor performance on both training and testing datasets.
    \item A model that overfits the training data has a low bias but high variance, resulting in an excellent performance on training examples while performing poorly on unseen data.
    \item Increasing the model complexity, such as by adding more parameters or hidden units, always improves a model's ability to generalize to unseen data.
    \item Using cross-validation guarantees that the selected hyperparameters will deliver optimal and flawless performance on any new unseen data.
    \item In support vector machines, choosing a very high value for the parameter \( C \) forces the model to aggressively minimize training errors, which can lead to a complex decision boundary and an increased risk of overfitting.
    \item For generative models based on Bayes’ rule, even with a large amount of training data, poor model assumptions or insufficient model complexity can result in underfitting.
    \item Overfitting occurs when a model learns not only the underlying pattern but also the noise inherent in the training data, leading to degraded performance on new, unseen examples.
    \item Achieving good generalization in classic machine learning models requires balancing the bias-variance trade-off through careful model selection, hyperparameter tuning, and appropriate regularization.
\end{enumerate}



\end{document}
